\section{Theoretical Foundations: Gaussian Mixture Models and EM Estimation}
\label{sec:gmm_theory}

This section establishes the mathematical foundations of Gaussian Mixture Models (GMM) as parametric density estimators for card fraud anomaly detection, derives the Expectation-Maximization (EM) algorithm, characterizes covariance topologies and parameter complexity, formalizes information-theoretic model selection via AIC and BIC, and scrutinizes the fundamental theoretical limitations of density-based anomaly scoring under the conditions of financial telemetry.

\subsection{Mixture Density and Latent Variable Formulation}
Let $\mathbf{x} \in \mathbb{R}^d$ denote a standardized $d$-dimensional transaction feature vector drawn from the legitimate transaction population. A Gaussian Mixture Model expresses the marginal probability density $p(\mathbf{x})$ as a convex combination of $K$ multivariate Gaussian component densities:
\begin{equation}
p(\mathbf{x} \mid \boldsymbol{\theta}) = \sum_{k=1}^K \pi_k \mathcal{N}(\mathbf{x} \mid \boldsymbol{\mu}_k, \boldsymbol{\Sigma}_k),
\label{eq:gmm_marginal}
\end{equation}
where $\boldsymbol{\theta} = \{\pi_k, \boldsymbol{\mu}_k, \boldsymbol{\Sigma}_k\}_{k=1}^K$ encapsulates all model parameters. The mixing coefficients $\pi_k = P(z_k = 1)$ satisfy the simplex constraints:
\begin{equation}
\sum_{k=1}^K \pi_k = 1, \quad \text{with } 0 \le \pi_k \le 1 \quad \forall k \in \{1, \dots, K\}.
\end{equation}
Here, $\mathbf{z} \in \{0, 1\}^K$ is a discrete 1-of-$K$ latent indicator vector such that $z_k = 1$ if transaction $\mathbf{x}$ was generated by the $k$-th Gaussian sub-population, with prior $P(z_k = 1) = \pi_k$. The conditional component density is the multivariate Gaussian:
\begin{equation}
\mathcal{N}(\mathbf{x} \mid \boldsymbol{\mu}_k, \boldsymbol{\Sigma}_k) = \frac{1}{(2\pi)^{d/2} |\boldsymbol{\Sigma}_k|^{1/2}} \exp\left( -\frac{1}{2} (\mathbf{x} - \boldsymbol{\mu}_k)^\top \boldsymbol{\Sigma}_k^{-1} (\mathbf{x} - \boldsymbol{\mu}_k) \right).
\label{eq:gaussian_density}
\end{equation}

By Bayes' theorem, the posterior probability (responsibility) that component $k$ was responsible for generating an observed observation $\mathbf{x}_n$ is given by:
\begin{equation}
\gamma_{nk} \equiv P(z_{nk} = 1 \mid \mathbf{x}_n, \boldsymbol{\theta}) = \frac{\pi_k \mathcal{N}(\mathbf{x}_n \mid \boldsymbol{\mu}_k, \boldsymbol{\Sigma}_k)}{\sum_{j=1}^K \pi_j \mathcal{N}(\mathbf{x}_n \mid \boldsymbol{\mu}_j, \boldsymbol{\Sigma}_j)}.
\label{eq:responsibility}
\end{equation}

\subsection{Maximum Likelihood and the EM Algorithm}
Given an independent and identically distributed (i.i.d.) training sample of legitimate transactions $\mathbf{X} = \{\mathbf{x}_1, \dots, \mathbf{x}_N\} \subset \mathbb{R}^d$, the log-likelihood function is:
\begin{equation}
\ln p(\mathbf{X} \mid \boldsymbol{\theta}) = \sum_{n=1}^N \ln \left( \sum_{k=1}^K \pi_k \mathcal{N}(\mathbf{x}_n \mid \boldsymbol{\mu}_k, \boldsymbol{\Sigma}_k) \right).
\label{eq:log_likelihood}
\end{equation}
Direct maximization of \eqref{eq:log_likelihood} with respect to $\boldsymbol{\theta}$ is analytically intractable because the summation over components appears inside the logarithm. The Expectation-Maximization (EM) algorithm \cite{dempster1977maximum} circumvents this via an auxiliary surrogate: the expected complete-data log-likelihood.

\subsubsection{E-Step (Expectation)}
Evaluate the responsibilities $\gamma_{nk}$ for every observation $n \in \{1, \dots, N\}$ and component $k \in \{1, \dots, K\}$ using the current parameter estimates $\boldsymbol{\theta}^{(t)}$ according to \eqref{eq:responsibility}.

\subsubsection{M-Step (Maximization)}
Re-estimate the parameters $\boldsymbol{\theta}^{(t+1)}$ by maximizing the expected complete-data log-likelihood $Q(\boldsymbol{\theta}, \boldsymbol{\theta}^{(t)})$:
\begin{align}
N_k &= \sum_{n=1}^N \gamma_{nk}, \label{eq:nk}\\
\pi_k^{(t+1)} &= \frac{N_k}{N}, \label{eq:m_pi}\\
\boldsymbol{\mu}_k^{(t+1)} &= \frac{1}{N_k} \sum_{n=1}^N \gamma_{nk} \mathbf{x}_n, \label{eq:m_mu}\\
\boldsymbol{\Sigma}_k^{(t+1)} &= \frac{1}{N_k} \sum_{n=1}^N \gamma_{nk} (\mathbf{x}_n - \boldsymbol{\mu}_k^{(t+1)})(\mathbf{x}_n - \boldsymbol{\mu}_k^{(t+1)})^\top. \label{eq:m_sigma}
\end{align}
For diagonal covariances ($\boldsymbol{\Sigma}_k = \text{diag}(\sigma_{k1}^2, \dots, \sigma_{kd}^2)$), \eqref{eq:m_sigma} simplifies to component-wise element-wise variances:
\begin{equation}
\sigma_{kj}^{2(t+1)} = \frac{1}{N_k} \sum_{n=1}^N \gamma_{nk} (x_{nj} - \mu_{kj}^{(t+1)})^2, \quad j \in \{1, \dots, d\}.
\end{equation}

\subsubsection{Monotonic Convergence Guarantee via Jensen's Inequality}
The EM algorithm guarantees that the observed-data log-likelihood is monotonically non-decreasing at each iteration:
\begin{equation}
\ln p(\mathbf{X} \mid \boldsymbol{\theta}^{(t+1)}) \ge \ln p(\mathbf{X} \mid \boldsymbol{\theta}^{(t)}).
\end{equation}
This property is derived using Jensen's inequality and the Evidence Lower Bound (ELBO) \cite{bishop2006pattern}. For any distribution $q(\mathbf{Z})$ over the latent indicators:
\begin{align}
\ln p(\mathbf{X} \mid \boldsymbol{\theta}) &= \sum_{\mathbf{Z}} q(\mathbf{Z}) \ln \left( \frac{p(\mathbf{X}, \mathbf{Z} \mid \boldsymbol{\theta})}{q(\mathbf{Z})} \right) + D_{\text{KL}}(q(\mathbf{Z}) \parallel p(\mathbf{Z} \mid \mathbf{X}, \boldsymbol{\theta})) \nonumber \\
&\ge \sum_{\mathbf{Z}} q(\mathbf{Z}) \ln \left( \frac{p(\mathbf{X}, \mathbf{Z} \mid \boldsymbol{\theta})}{q(\mathbf{Z})} \right) \equiv \mathcal{L}(q, \boldsymbol{\theta}),
\end{align}
since the Kullback-Leibler divergence $D_{\text{KL}} \ge 0$. In the E-step, setting $q(\mathbf{Z}) = p(\mathbf{Z} \mid \mathbf{X}, \boldsymbol{\theta}^{(t)})$ forces $D_{\text{KL}} = 0$, rendering the lower bound tight: $\mathcal{L}(q^{(t)}, \boldsymbol{\theta}^{(t)}) = \ln p(\mathbf{X} \mid \boldsymbol{\theta}^{(t)})$. In the M-step, maximizing $\mathcal{L}(q^{(t)}, \boldsymbol{\theta})$ over $\boldsymbol{\theta}$ yields:
\begin{equation}
\ln p(\mathbf{X} \mid \boldsymbol{\theta}^{(t+1)}) \ge \mathcal{L}(q^{(t)}, \boldsymbol{\theta}^{(t+1)}) \ge \mathcal{L}(q^{(t)}, \boldsymbol{\theta}^{(t)}) = \ln p(\mathbf{X} \mid \boldsymbol{\theta}^{(t)}).
\end{equation}

\subsubsection{Local Optimality and Multi-Start Initialization}
Although monotonicity prevents divergence, the log-likelihood surface of a mixture model is non-convex and exhibits multimodal ridges and saddle points \cite{mclachlan2000finite}. Consequently, EM converges strictly to a **local maximum** (or stationary point) dependent upon the initial parameter seed $\boldsymbol{\theta}^{(0)}$. To protect against suboptimal basins of attraction, multiple independent initializations ($n_{\text{init}} \ge 3$) seeded via $k$-means clustering are essential.

\subsection{The Singularity Problem and Covariance Regularization}
A critical pathological breakdown of Maximum Likelihood in mixture models occurs when a Gaussian component collapses onto an isolated sample $\mathbf{x}_m$. If $\boldsymbol{\mu}_k \to \mathbf{x}_m$, the squared distance $(\mathbf{x}_m - \boldsymbol{\mu}_k)^\top \boldsymbol{\Sigma}_k^{-1} (\mathbf{x}_m - \boldsymbol{\mu}_k) \to 0$. As the variance of this component shrinks toward zero, the determinant $|\boldsymbol{\Sigma}_k| \to 0$, causing the component density and the total likelihood to explode toward infinity:
\begin{equation}
\lim_{|\boldsymbol{\Sigma}_k| \to 0} \mathcal{N}(\mathbf{x}_m \mid \boldsymbol{\mu}_k, \boldsymbol{\Sigma}_k) = +\infty \implies \ln p(\mathbf{X} \mid \boldsymbol{\theta}) \to +\infty.
\end{equation}
This singularity is an artifact of unconstrained optimization rather than a meaningful physical density. To guarantee positive definiteness and invertibility, covariance regularization is introduced by augmenting the diagonal elements:
\begin{equation}
\boldsymbol{\Sigma}_k^{\text{reg}} = \boldsymbol{\Sigma}_k + \epsilon_{\text{cov}} \mathbf{I}_d,
\label{eq:reg_covar}
\end{equation}
where $\epsilon_{\text{cov}} \in [10^{-6}, 10^{-2}]$ establishes a strictly positive lower bound on eigenvalues: $\lambda_{\min}(\boldsymbol{\Sigma}_k^{\text{reg}}) \ge \epsilon_{\text{cov}} > 0$.

\subsection{Covariance Topologies and Free Parameter Complexity}
The choice of covariance matrix structure imposes geometric assumptions on the component shapes while dramatically impacting parameter complexity $p$. For a feature space of dimension $d$ and component count $K$:
\begin{enumerate}
    \item \textbf{Full Covariance (`full'):} Arbitrary ellipsoids with distinct orientations and eccentricities. Parameters: $K d$ (means) $+ K \frac{d(d+1)}{2}$ (symmetric covariance elements) $+ (K - 1)$ (mixing weights):
    \begin{equation}
    p_{\text{full}} = K d + \frac{K d(d+1)}{2} + (K - 1) = \mathcal{O}(K d^2).
    \end{equation}
    \item \textbf{Tied Covariance (`tied'):} Ellipsoids sharing a single common orientation $\boldsymbol{\Sigma}_1 = \dots = \boldsymbol{\Sigma}_K = \boldsymbol{\Sigma}$:
    \begin{equation}
    p_{\text{tied}} = K d + \frac{d(d+1)}{2} + (K - 1) = \mathcal{O}(K d + d^2).
    \end{equation}
    \item \textbf{Diagonal Covariance (`diag'):} Axis-aligned ellipsoids assuming conditional feature independence given the component ($\boldsymbol{\Sigma}_k = \text{diag}(\sigma_{k1}^2, \dots, \sigma_{kd}^2)$):
    \begin{equation}
    p_{\text{diag}} = K d + K d + (K - 1) = 2Kd + (K - 1) = \mathcal{O}(Kd).
    \end{equation}
    \item \textbf{Spherical Covariance (`spherical'):} Isotropic hyperspheres ($\boldsymbol{\Sigma}_k = \sigma_k^2 \mathbf{I}_d$):
    \begin{equation}
    p_{\text{spherical}} = K d + K + (K - 1) = K(d + 2) - 1 = \mathcal{O}(Kd).
    \end{equation}
\end{enumerate}

\begin{table}[h]
\centering
\caption{GMM Parameter Complexity for $d=31$ Engineered Features.}
\label{tab:gmm_param_complexity}
\footnotesize
\begin{tabular}{lrrrr}
\hline
\textbf{Topology} & \textbf{Formula $p$} & \textbf{$K=1$} & \textbf{$K=5$} & \textbf{$K=16$} \\
\hline
Full & $Kd + Kd(d+1)/2 + (K-1)$ & $527$ & $2\,639$ & $8\,447$ \\
Tied & $Kd + d(d+1)/2 + (K-1)$ & $527$ & $651$ & $992$ \\
Diag & $2Kd + (K-1)$ & $62$ & $314$ & $1\,007$ \\
Spherical & $K(d+2) - 1$ & $32$ & $164$ & $527$ \\
\hline
\end{tabular}
\end{table}

Plugging $d=31$ into Table~\ref{tab:gmm_param_complexity} demonstrates why `full' covariance is perilous: at $K=16$, full covariance requires estimating $8\,447$ continuous parameters from finite iterations, whereas diagonal covariance requires only $1\,007$ parameters. Because V1--V28 are PCA orthogonal projections ($r \approx 0$), diagonal covariance is mathematically consistent and provides massive variance reduction.

\subsection{Information-Theoretic Model Selection: AIC vs. BIC}
Selecting the optimal component count $K$ and covariance structure requires balancing training likelihood against over-parameterization. Let $\hat{L} = \prod_{n=1}^N p(\mathbf{x}_n \mid \hat{\boldsymbol{\theta}})$ denote the maximized likelihood on a sample of size $N$. We evaluate two criteria:
\begin{align}
\text{AIC} &= 2p - 2 \ln \hat{L}, \label{eq:aic} \\
\text{BIC} &= p \ln N - 2 \ln \hat{L}. \label{eq:bic}
\end{align}

\subsubsection{Theoretical Distinctions and Asymptotic Behavior}
\begin{itemize}
    \item \textbf{Penalty Magnitude:} While AIC penalizes complexity with a constant factor of $2p$, BIC penalizes complexity by $p \ln N$. For our training sample ($N = 50\,000$), $\ln(50\,000) \approx 10.82 \gg 2$. Thus, BIC penalizes additional components more than five times as severely as AIC.
    \item \textbf{Consistency vs. Efficiency:} BIC is **dimensionally consistent** \cite{schwarz1978estimating}: as $N \to \infty$, if the true data-generating process is a finite Gaussian mixture, the probability of selecting the true model converges to $1$. In contrast, AIC is **minimax-rate optimal for prediction** \cite{akaike1974new}: it asymptotically minimizes out-of-sample Kullback-Leibler divergence but tends to overfit by selecting unnecessarily large $K$.
    \item \textbf{The Elbow and Parsimony:} In high-dimensional spaces, AIC monotonically decreases as $K$ increases, failing to identify an elbow. BIC, however, reaches a distinct minimum and rises sharply as variance penalties dominate likelihood gains.
\end{itemize}

\subsection{Anomaly Scoring Convention and Foundational Limitations}
\subsubsection{Negative Log-Likelihood (NLL) Scoring}
Under the semi-supervised anomaly detection paradigm, the anomaly score $s(\mathbf{x})$ is defined as the negative log-marginal density:
\begin{equation}
s(\mathbf{x}) = -\ln p(\mathbf{x} \mid \hat{\boldsymbol{\theta}}) = -\ln \left( \sum_{k=1}^K \hat{\pi}_k \mathcal{N}(\mathbf{x} \mid \hat{\boldsymbol{\mu}}_k, \hat{\boldsymbol{\Sigma}}_k) \right).
\label{eq:nll_score}
\end{equation}
This satisfies the required score convention: transactions situated in dense, high-probability regions of the legitimate training manifold have high $p(\mathbf{x})$ and receive small scores $s(\mathbf{x}) \approx 0$; transactions falling into low-density sparse margins receive high scores $s(\mathbf{x}) \gg 0$.

\subsubsection{The Fundamental Limitation: Density $\neq$ Fraud}
While mathematically elegant, density-based scoring possesses four intrinsic theoretical limitations in financial fraud detection:
\begin{enumerate}
    \item \textbf{Unusualness is not Malice:} Low density implies only that a transaction is statistically infrequent (e.g. an affluent cardholder buying a $\$15\,000$ luxury watch at noon). High NLL does not imply fraudulent intent, generating unavoidable false positives.
    \item \textbf{Subtle Fraud Infiltration:} Sophisticated fraudsters intentionally mimic legitimate spending behaviors (micro-authorizations under $\$5.00$ at everyday grocers). These frauds inhabit the high-density interior mass where $s(\mathbf{x})$ is minimal, eluding density detection entirely.
    \item \textbf{Gaussian Tail Decay vs. Leptokurtosis:} A Gaussian density tails off exponentially ($\propto \exp(-\frac{1}{2} z^2)$). As established in Section~\ref{sec:dataset_eda}, real financial features exhibit massive excess kurtosis ($\text{kurt} > 100$). Fitting thin-tailed Gaussians to fat-tailed data forces legitimate extreme spenders to receive astronomical anomaly scores.
    \item \textbf{Curse of Dimensionality and Distance Concentration:} In dimension $d$, the volume of space grows as $r^d$, and Gaussian mass concentrates entirely in a thin spherical shell (the "soap bubble" effect) away from the mode. Density estimation becomes exponentially sample-inefficient as $d$ grows \cite{aggarwal2017outlier}.
\end{enumerate}
